\def\stat{borisov3}





\def\tit{СТРАТЕГИЯ ИССЛЕДОВАНИЙ И~РАЗРАБОТОК\\ В~ОБЛАСТИ


ИСКУССТВЕННОГО ИНТЕЛЛЕКТА~III:\\ ДОКТРИНА


ГОСУДАРСТВЕННОЙ ПОДДЕРЖКИ~США}








\def\titkol{Стратегия исследований и~разработок в~области


искусственного интеллекта~III} %: Доктрина государственной поддержки США}





\def\aut{А.\,В.~Борисов$^1$, А.\,В.~Босов$^2$, Д.\,В.~Жуков$^3$}





\def\autkol{А.\,В.~Борисов, А.\,В.~Босов, Д.\,В.~Жуков}





\titel{\tit}{\aut}{\autkol}{\titkol}





\index{Борисов А.\,В.}


\index{Босов А.\,В.}


\index{Жуков Д.\,В.}


\index{Borisov A.\,V.}


\index{Bosov A.\,V.}


\index{Zhukov D.\,V.}





%{\renewcommand{\thefootnote}{\fnsymbol{footnote}}


%\footnotetext[1]


%{Работа частично поддержана РФФИ (проект 18-29-03081).}}





\renewcommand{\thefootnote}{\arabic{footnote}}


\footnotetext[1]{Федеральный исследовательский центр <<Информатика и~управ\-ле\-ние>>


Российской академии наук, \mbox{ABorisov@ipiran.ru}}


\footnotetext[2]{Федеральный исследовательский центр <<Информатика и~управ\-ле\-ние>>


Российской академии наук, \mbox{AVBosov@ipiran.ru}}


\footnotetext[3]{Федеральный исследовательский центр <<Информатика и~управ\-ле\-ние>>


Российской академии наук, \mbox{DZhukov@ipiran.ru}}





 \label{st\stat}





 \thispagestyle{headings}





 \vspace*{-12pt}





   \Abst{Статья продолжает цикл работ, посвященных анализу влияния


государственного управления на эффективность проведения исследований и~разработок


в~области искусственного интеллекта (AI~R\&D, Artificial Intelligence Research and Development). 


В~третьей час\-ти цик\-ла рассматривается


влияние государства на AI~R\&D на примере США. Дано краткое описание стратегического


документа США в~области AI~R\&D, вклю\-чая компактное изложение его целей и~задач, а~так\-же


принципов реализации. Помимо этого проанализирован смеж\-ный документ Министерства


обороны США. Пред\-став\-лен классификационный анализ направлений AI~R\&D в~области


обороны и~безопас\-ности, проводимых основными специализированными на\-уч\-но-ис\-сле\-до\-ва\-тель\-ски\-ми 


организациями США.}





   \KW{искусственный интеллект; Министерство обороны США (DoD); Управление


перспективных исследовательских проектов Министерства обороны США (DARPA); Агентство


передовых исследований в~сфере разведки (IARPA)}





   \DOI{14357\!/\!08696527210410}





   \vskip 10pt plus 9pt minus 6pt





    \thispagestyle{headings}





    \vspace*{-12pt}





   \section{Введение}





  Статья представляет третью часть цик\-ла работ, посвященного анализу


государственного влияния на сферу исследований и~разработок в~об\-ласти


искусственного интеллекта. Предыдущие части~\cite{1-bor, 2-bor}


были посвящены хронологии развития данной области в~СССР


и~Российской Федерации и~сравнительному анализу наукометрических


показателей публикационной активности по AI~R\&D за последние 20~лет


в~России и~мире. Эти обзоры сформировали пред\-став\-ле\-ние о~<<начальных


условиях>> реализации~<<Национальной стратегии развития


искусственного интеллекта на период до~2030~года>>~\cite{3-bor}.





  Наше государство было не первым, кто разработал и~приступил


к~реализации стратегии подобного рода. Аналогичные документы были


приняты к~выполнению в~2016~г.\ в~США и~в~2017~г.\ в~КНР. Цель


данной части цикла~--- анализ государственного влияния в~об\-ласти


AI~R\&D в~США. Соединенные Штаты Америки являются не только


сверхдержавой, родоначальницей средств вычислительной техники, но


и~инициатором развития отрасли AI~R\&D в~целом.





  Работа организована следующим образом. В~разд.~2 дано краткое


описание стратегического документа США~\cite{4-bor}, а~также заявленных


в~нем целей, задач и~принципов их реализации. К~одной из исключительных


государственных прерогатив относится функция обеспечения обороны,


общественной и~государственной безопас\-ности. Раздел~3 статьи посвящен


анализу программных документов Министерства обороны США в~час\-ти


AI~R\&D. Пред\-став\-ле\-ны основные на\-прав\-ле\-ния применения технологий


искусственного интеллекта в~об\-ласти обороны и~безопас\-ности, а~так\-же


названы основные государственные специализированные


на\-уч\-но-ис\-сле\-до\-ва\-тель\-ские организации, фор\-ми\-ру\-ющие спектр AI~R\&D:


Управление перспективных исследовательских проектов Министерства


обороны США (\mbox{DARPA}, Defence Advanced Research Projects Agency) 


и~Агентство передовых исследований в~сфере


разведки (\mbox{IARPA}, Intelligence Advanced Research Projects Activity). 


На основе доступных источников проведен анализ


на\-прав\-ле\-ний прикладных исследований данных организаций в~об\-ласти


искусственного интеллекта. Раздел~4 посвящен внедрению приложений,


ис\-поль\-зу\-ющих искусственный интеллект, в~правоохранительных органах


США: ФБР и~полиции. Заключительные выводы по данной час\-ти цик\-ла


пред\-став\-ле\-ны в~разд.~5.





\section{Государственный подход к~исследованиям и~разработкам в~области~искусственного


интеллекта в~США} %2





  Исследования в~об\-ласти искусственного интеллекта проводятся в~США


дав\-но. Американское правительство выделяло и~выделяет на эти цели


весьма большие средства. Например, в~2015~г.\ общий объем затрат США на


AI~R\&D составил примерно 1,1~млрд долларов, что с~учетом инфляции


сейчас эквивалентно 80~млрд рублей. При этом размер частных инвестиций


в~сферу AI~R\&D точно назвать сложно, однако он, без сомнения, больше


государственного. Известно, что объем венчурного (рискового) капитала,


направленного на стартапы в~об\-ласти AI~R\&D, в~период с~2013


по~2014~гг.\ вырос в~4~раза. Для сравнения: на проведение всех


фундаментальных исследований в~2020~г.\ (включая РАН, государственные


исследовательские центры, НИИ, Сколково и~др.) до начала пандемии


COVID-19 планировалось выделить 150~млрд рублей.





  Понимание стратегической важности AI~R\&D в~США нашло свое


отражение в~разработанной Национальным советом по науке и~технике


США (NSTC) в~2016~г.\ документа~<<Национальный стратегический план


по исследованиям и~разработке искусственного интеллекта>>~\cite{4-bor}


(далее~--- Стратегия). В~нем выделены на\-прав\-ле\-ния AI~R\&D, которые


будут получать государственное финансирование. Документ не


фокусировался на AI~R\&D исключительно в~военных целях. Основной


целью данных исследований было названо получение новых знаний


и~технологий в~об\-ласти AI~R\&D. Для ее достижения были определены


сле\-ду\-ющие на\-прав\-ле\-ния и~ключевые принципы деятельности (стратегии,


в~терминах документа):





\noindent


\begin{enumerate}[(1)]


\item %1


\textbf{долговременные инвестиции в~AI~R\&D:}


горизонт ожидания результатов~--- 5--10~лет; <<высокий риск\,--\,боль\-шие выгоды>>; развитие


методологий, ориентированных на данные, для получения новых знаний; улучшение


возможностей восприятия у~сис\-тем~AI; исследование тео\-ре\-ти\-че\-ских возможностей


и~ограничений AI; проведение исследований в~об\-ласти AI общего назначения; создание


масштабируемых систем AI; поддержка исследований AI, подобного человеческому; создание


более совершенных и~на\-деж\-ных роботов; развитие аппаратной базы для улучшенного AI;


разработка сис\-тем AI для модернизированной аппаратной базы;


\item %2


\textbf{разработка эффективных методов взаимодействия человека и~AI:}


поиск новых алгоритмов взаимодействия~<<че\-ло\-век--AI>>; использование AI для развития


человеческих способностей; развитие технологии визуализации и~интерфейсов~<<че\-ло\-век--AI>>;


создание более эффективных сис\-тем лингвистической обработки;


\item %3


\textbf{выявление и~решение этических, законодательных и~социальных


проб\-лем использования AI:}


повышение уровня спра\-вед\-ли\-вости, про\-зрач\-ности и~учета; по\-стро\-ение систем AI, подчиненных


этическим нормам; разработка архитектуры этического AI;


\item %4


\textbf{обеспечение безопасности использования систем AI:}


повышение объяснимости и~прозрачности; увеличение доверия; улучшение качества


верификации и~под\-тверж\-де\-ния; увеличение за\-щи\-щен\-ности против атак различной природы;


достижение долгосрочной безопас\-ности сис\-тем AI и~выравнивание качества;


\item %5


\textbf{разработка открытых массивов данных для обучения


и~тестирования систем AI:}


создание и~пуб\-ли\-ка\-ция широкого спектра массивов данных для нужд разработчиков AI; создание


ресурсов для тренировок и~тес\-тов с~учетов коммерческих и~общественных интересов;


разработка биб\-лио\-тек и~пакетов с~открытым кодом;


\item %6


\textbf{метрологическое обеспечение технологий искусственного


интеллекта путем создания стандартов и~критериев:}


создание широкого спектра стандартов по AI; разработка технологических тес\-тов; повышение


до\-ступ\-ности тес\-то\-вых стендов для AI; вовлечение сообщества разработчиков и~исследователей AI


в~процесс создания стандартов и~тес\-тов;


\item %7


\textbf{изучение и~учет потребностей национальных разработчиков AI.}


\end{enumerate}





В~качестве своих со\-став\-ля\-ющих стратегия включает сле\-ду\-ющие


стратегические планы:


  \begin{itemize}


\item Федеральный план исследований и~разработок в~об\-ласти Big


Data~\cite{5-bor};


\item Федеральный план исследований и~разработок в~об\-ласти


кибербезопасности~\cite{6-bor};


\item Национальную стратегию исследования защиты персональных


данных~\cite{7-bor};


\item Национальный инициативный план по исследованиям в~об\-ласти


нанотехнологий~\cite{8-bor};


\item Национальную компьютерную стратегическую инициативу~\cite{9-bor};


\item Инициативу по исследованию мозга с~по\-мощью инновационных


нейротехнологий~\cite{10-bor};


\item Национальную робототехническую инициативу~\cite{11-bor}.


  \end{itemize}





  Согласно разработанной Стратегии, США планируют достичь следующих


преимуществ путем внед\-ре\-ния технологий искусственного интеллекта:


  \begin{itemize}


\item увеличить экономическое бла\-го\-со\-сто\-яние (развитие производства,


логистики, финансов, транспорта, сельского хозяйства, торговли,


коммуникаций, науки и~технологий);


\item улучшить возможности обуче\-ния и~качество жизни (образование,


медицина, юриспруденция, персонализированные сервисы);


\item укрепить государственную и~общественную безопас\-ность.


  \end{itemize}





  Летом 2019~г.\ по первым результатам выполнения Стратегии был


опубликован отчет~\cite{12-bor}, содержащий ис\-чер\-пы\-ва\-ющий анализ


результатов трех лет ее выполнения. По результатам первого этапа


реализации Стратегии была скорректирована организационная структура для


ее реализации. Она пред\-став\-ле\-на в~табл.~1.








  %tabl 1 ---> .eps








  \begin{table}\small %tabl1


  \begin{center}


  \Caption{Скорректированная версия организационной структуры реализации Стратегии}


\label{t1}


  \vspace*{2ex}





\begin{center}


\mbox{%


%\epsfxsize=90.622mm 


%\epsfbox{bo3-1.eps} %гор


\epsfxsize=128mm 


\epsfbox{bo3-1t.eps} %верт


}


\end{center}


%\vspace*{-9pt}


  \end{center}


  \end{table}





  Результатом анализа AI~R\&D стало определение еще одного, восьмого,


на\-прав\-ле\-ния в~этой об\-ласти:


\begin{enumerate}[(1)]


\setcounter{enumi}{7}


\item 


\textbf{расширение го\-су\-дар\-ст\-вен\-но-част\-но\-го партнерства для ускорения развития.}


  Выделение этого направления связано с~осознанием важности тезиса


Стратегии об~<<устойчивом финансировании AI~R\&D в~сотрудничестве


с~академическими, производственными международными партнерами


и~союзниками и~другими негосударственными юридическими лицами для


достижения технологических прорывов в~об\-ласти AI и~смежных


технологиях для их быстрейшей трансформации в~реальные средства


укрепления национальной и~экономической безопас\-ности США>>.


Част\-но-го\-су\-дар\-ст\-вен\-ное партнерство позволит эффективно перераспределять


ресурсы и~знания для скорейшего достижения результатов.


  \end{enumerate}





\section{Направления исследований и~разработок в~области искусственного интеллекта для~обороны и~безопасности} %3





  Министерство обороны США (DoD, Department of Defence) в~рамках реализации Национальной


оборонной инициативы США 2018 г. разработало и~опубликовало документ~\cite{13-bor}, 


содержащий краткие сведения о~стратегии страны по


использованию технологий искусственного интеллекта в~военных целях.


Подразделением DoD, ответственным за выполнение этой стратегии, назван


Объединенный центр искусственного интеллекта (JAIC, Joint Artificial Intelligence Center). 


В~качестве


организаций, ведущих работы по созданию и~адап\-та\-ции разработок


в~области AI для нужд DoD, указаны Управ\-ле\-ние перспективных


исследовательских проектов Министерства обороны США (DARPA),


Агентство передовых исследований в~сфере разведки (IARPA),


национальные исследовательские (в~исходном тексте~--- \textit{academic})


центры, частные компании, международные объединения.





  В~документе указаны следующие общие направления применения


искусственного интеллекта для достижения стратегического военного


преимущества.


\begin{description}


\item[\,]


  \textbf{Улучшение ситуационной осведомленности и~принятия решений}.


Технологии AI, примененные для решения задач вос\-при\-я\-тия\!/\!осмыс\-ле\-ния,


таких как анализ изображений, могут извлекать полезную информацию из


сырых начальных данных и~тем самым повышать ситуационную


осведомленность командного состава. Искусственный интеллект может генерировать и~помогать


руководству исследовать оперативные сценарии, чтобы оно впоследствии


могло осуществить оптимальный выбор для одновременного достижения


по\-став\-лен\-ной цели и~минимальных рисков как для вооруженных сил, так


и~для гражданского населения.


\item[\,]


  \textbf{Повышение безопасности эксплуатации оборудования}. Искусственный интеллект обладает


серьезным потенциалом повышения безопас\-ности эксплуатации военной


техники (самолетов, кораблей, наземных транспортных и~боевых средств)


в~слож\-ных быст\-ро\-ме\-ня\-ющих\-ся ситуациях, предуп\-реж\-дая\linebreak операторов


о~скрытых опасностях.


\item[\,]


  \textbf{Внедрение упреждающего обслуживания и~снабжения}~---


использование AI для прогнозирования выхода из строя критических частей,


автоматической диагностики и~планирования технического обслуживания на


основе статистических данных и~со\-сто\-яния оборудования. Подобная же


технология будет использоваться для упреж\-да\-юще\-го управ\-ле\-ния поставками


запасных частей и~оптимизации уровня запасов. Эти преимущества


обеспечат поддержку запасов на тре\-бу\-емом уровне, а~так\-же помогут


в~устранении неполадок и~сделают возможным более быст\-рое


развертывание средств и~адап\-та\-цию затрат при снижении издержек.


\item[\,]


  \textbf{Оптимизация деловых процессов}. Искусственный интеллект будет использоваться с~\mbox{целью}


сокращения времени, затрачиваемого на ручные, повторяющиеся и~час\-тые


задачи. Оставляя людям задачи надзора за выполнением


автоматизированных функций, AI имеет потенциал к~сокращению чис\-ла


и~стоимости ошибок, увеличению выхода и~степени приспособления, для


направления ресурсов DoD на выполнение более приоритетных задач


и~достижение важнейших целей.


\end{description}





  Ниже представлена информация о~проектах в~области AI~R\&D,


проводимых~<<головными>> специализированными исследовательскими


организациями, упомянутыми в~\cite{13-bor},~--- DARPA~\cite{14-bor}


и~\mbox{IARPA}~\cite{15-bor}.





  Управление DAPRA объявило в~сентябре 2018~г.\ о~многолетних инвестициях


в~размере более 2~млрд долл.\ США в~новые и~существующие


программы, на\-зы\-ва\-емые кампанией~<<AI Next>>. Ключевые составляющие


кампании предполагают автоматизацию критически важных


биз\-нес-про\-цес\-сов Министерства обороны, таких как:


  \begin{itemize}


\item проверка соответствия требованиям безопас\-ности или аккредитация


программных систем для оперативного развертывания;


\item повышение на\-деж\-ности сис\-тем искусственного интеллекта;


\item повышение безопас\-ности и~отказоустойчивости машинного


обучения и~технологий искусственного интеллекта;


\item снижение зависимости от нехватки данных и~производительности.


  \end{itemize}





 Управление перспективных исследовательских проектов Министерства обороны


  США создаст мощные возможности для DoD, опираясь на сле\-ду\-ющие


принципы.


  \begin{enumerate}[1.]





\item %1


\textbf{Обеспечение новых возможностей.}


Технологии AI применяются регулярно для реализации проектов \mbox{DARPA}, включая более 60


существующих программ, связанных с~анализом сложных кибератак в~реальном времени,


обнаружением мошеннических изображений, построением динамических цепочек убийств, для


ведения войн во всех областях, технологиях человеческого языка, мультимодальном


автоматическом распознавании целей и~др. Управ\-ле\-ние \mbox{DARPA} будет продвигать технологии AI для


обеспечения автоматизации критически важных биз\-нес-про\-цес\-сов DoD. Один из таких


процессов~--- аккредитация программных систем перед эксплуатационным развертыванием.


Автоматизация этого процесса аккредитации с~использованием известных AI и~других


технологий теперь представляется возможной.


\item %2


\textbf{Обеспечение надежного искусственного интеллекта.}


Технологии AI продемонстрировали большое значение для таких разнообразных задач, как анализ


спутниковых изображений, предупреждение о~кибератаках, логистика и~др.\ В~то же время


режимы от\-ка\-за\!/\!ава\-рий технологий AI плохо изучены. Управ\-ле\-ние


\mbox{DARPA} работает над решением этой


проблемы, уделяя особое внимание исследованиям и~разработкам~--- как аналитическим, так


и~эмпирическим.


\item %3


\textbf{Учет фактора состязательности при разработке искусственного интеллекта.}


Самый мощный инструмент AI на сегодня~--- это машинное обучение (ML,


machine learning). Сис\-те\-мы ML могут


быть легко обмануты изменениями, которые никогда не обманут человека (так на\-зы\-ва\-емые \textit{adverse


samples}). Данные, используемые для обучения таких сис\-тем, могут быть по\-вреж\-де\-ны, и~само


программное обеспечение (ПО) уязвимо для кибератак. Задачи обеспечения устой\-чи\-вости ПО


и~защиты банков данных, ис\-поль\-зу\-емых для обуче\-ния, планируется решать совместно


и~оперативно.


\item %4


\textbf{Создание высокопроизводительного искусственного интеллекта.}


Повышение производительности компьютеров за последнее десятилетие поз\-во\-ли\-ло добиться


успеха в~машинном обучении в~сочетании с~большими наборами данных и~биб\-лио\-те\-ка\-ми


ПО. Повышение производительности при снижении энергопотребления


необходимо для развертывания цент\-раль\-но\-го и~тактического звеньев. 


Управ\-ле\-ние \mbox{DARPA}


продемонстрировало аналоговую обработку алгоритмов AI с~ускорением в~1000~раз


и~эффективностью энергопотребления в~1000~раз по сравнению с~современными цифровыми


процессорами, а~так\-же исследует ап\-па\-рат\-ные разработки для AI. \mbox{DARPA} также прикладывает


усилия для снижения требований к~маркированным данным обуче\-ния (качеству и~объемам).


\item %5


\textbf{Создание искусственного интеллекта следующего поколения.}


Алгоритмы машинного обучения, обеспечивающие распознавание лиц, и~автомобили


с~автоматическим вож\-де\-ни\-ем были изобретены более 20~лет назад. 


Управ\-ле\-ние \mbox{DARPA} заняло ведущее


место в~новаторских исследованиях по разработке следующего поколения алгоритмов AI,


которые превратят компьютеры из инструментов в~партнеров по решению проб\-лем.


Исследование \mbox{DARPA} направлено на то, чтобы дать возможность рассуждать (делать


умозаключения).


  \end{enumerate}








  Агентство передовых исследований в~сфере разведки (\mbox{IARPA}) было


создано в~2006~г.\ в~целях проведения междисциплинарных исследований,


определения новых возможностей и~создания прорывных технологий


в~области разведки. Задачи \mbox{IARPA} должны решаться, основываясь на


технических и~эксплуатационных знаниях, которыми обладают органы


разведки. Таким образом, миссия \mbox{IARPA} заключается в~прогнозировании


долгосрочных по\-треб\-но\-стей разведывательного сообщества


и~предоставлении ему технических и~исследовательских воз\-мож\-но\-стей.





  В табл.~2 представлена классификация направлений исследований \mbox{DARPA}


и~\mbox{IARPA}, в~той или иной степени ба\-зи\-ру\-ющих\-ся на AI~R\&D или


касающихся AI~R\&D. Следует отметить, что доля подобных проектов


у~\mbox{IARPA} от общего объема проводимых исследований со\-став\-ля\-ет


45\%--55\%, в~\mbox{DARPA} этот показатель ниже. В~данную классификацию не


включены работы, связанные с~различными беспилотными аппаратами,


а~так\-же AI~R\&D в~об\-ласти химии, материаловедения и~пр.








  \begin{table}\small %tabl2


  \begin{center}


  \Caption{Классификация AI~R\&D, проводимых DARPA и~IARPA }


\label{t2}


  \vspace*{2ex}





\tabcolsep=2pt


\begin{tabular}{|l|p{221.3pt}|c|c|}


  \hline


\multicolumn{1}{|c|}{Отрасль AI}&\multicolumn{1}{c|}{Направление исследований}&


\hspace*{2pt}{\raisebox{-17pt}{\rotatebox{90}{\tabcolsep=2pt\begin{tabular}{c} Проекты~\\ DARPA \end{tabular}}}}&


\hspace*{2pt}{\raisebox{-17pt}{\rotatebox{90}{\tabcolsep=2pt\begin{tabular}{c} Проекты~\\ IARPA \end{tabular}}}}\\


\hline


\multicolumn{1}{|l|}{\hphantom{1.~Системы обработ-}}&&&\\[-10pt]


\multicolumn{1}{|l|}{\raisebox{-80pt}[0pt][0pt]{\tabcolsep=0pt\begin{tabular}{l} 


1.~Системы обработ-\\ ки текстовой\!/\!рече-\\ вой информации \end{tabular}}}&


%\tabcolsep=0pt\begin{tabular}{l}


--~Глубокая обработка текста, выявление неочевидного смысла, выделение оперативно значимой


информации\newline


--~Интеллектуальная обработка малоресурсных языковых запросов с~целью повышения ситуационной


осведомленности\newline


--~Системы распознавания речи на любом человеческом языке\newline


--~Изучение использования метафор различными культурами, чтобы понять культурные нормы различных


народов\newline


--~Раскрытие социальных целей членов групп по их сообщениям\newline


--~Разработка методов извлечения детализируемой семантической информации с~акцентом на событиях


в~форме~<<кто--что--кому--когда-то>>


%\end{tabular}


&\hspace*{2pt}{\raisebox{-115pt}{\rotatebox{90}{DEFT, LORELEI}}}


&\hspace*{2pt}{\raisebox{-130pt}{\rotatebox{90}{\tabcolsep=2pt\begin{tabular}{c} Babel, Metaphor, SCIL,\\ BETTER, MATERIAL\end{tabular}}}}\\


\hline


\multicolumn{1}{|l|}{\raisebox{-72pt}[0pt][0pt]{\tabcolsep=0pt\begin{tabular}{l} 2.~Системы анализа\\ и~обработки\\


изображений\end{tabular}}}&


%\tabcolsep=0pt\begin{tabular}{l}


--~Разработка средств оценивания, манипуляций и~подделки цифровых изображений\newline


--~Аналитическая обработка видео- и~аудиоинформации\newline


--~Разработка системы распознавания спутниковых фотографий.\newline


--~Разработка технологии автоматической генерации точных трехмерных моделей объектов с~реальными


физическими свойствами из множества источников данных изображений\newline


--~Разработка ПО автоматического обнаружения движения для многокамерной среды потокового


видео\newline


--~Повышение производительности систем распознавания лиц\newline


--~Системы обнаружения антропогенной деятельности\hspace*{-.7pt}


%\end{tabular}


&\hspace*{2pt}{\raisebox{-90pt}{\rotatebox{90}{MefiFor}}}


&\hspace*{2pt}{\raisebox{-142pt}{\rotatebox{90}{\tabcolsep=0pt\begin{tabular}{c}Aladdin Video, Amon-Hen,\\ 


CORE3D, DIVA, Janus, SMART\end{tabular}}}}\\


\hline


\multicolumn{1}{|l|}{\raisebox{-12pt}[0pt][0pt]{\tabcolsep=0pt\begin{tabular}{l} 3.~Создание\\ биометрических\\


систем\end{tabular}}}&


%\tabcolsep=0pt\begin{tabular}{l}


--~Разработка высокоточных систем сопоставления, получения биометрических сигнатур по данным


низкого качества


%\end{tabular}


&&\hspace*{2pt}{\raisebox{-20pt}{\rotatebox{90}{BEST}}}\\


\hline


\multicolumn{1}{|l|}{\raisebox{-28pt}[0pt][0pt]{\tabcolsep=0pt\begin{tabular}{l} 4.~Системы обработ-\\ ки специальных\\ сигналов


и~данных\end{tabular}}}&


%\tabcolsep=0pt\begin{tabular}{l}


--~Внедрение ML-методов для обработки сигналов радиочастотного спектра\newline


--~Системы глубокой аналитической обработки данных геолокации\newline


--~Разработка методов непрерывного анализа данных радиоэлектронной разведки


%\end{tabular}


&\hspace*{2pt}{\raisebox{-40pt}{\rotatebox{90}{RFMLS}}}


&\hspace*{2pt}{\raisebox{-40pt}{\rotatebox{90}{\tabcolsep=0pt\begin{tabular}{c} Finder,\\ Mercury\end{tabular}}}}\\


\hline


\multicolumn{4}{r}{\footnotesize\textit{Продолжение табл.~$2$ на с.}~\pageref{t2-1}}\\


  \end{tabular}


  \end{center}


  \end{table}





\addtocounter{table}{-1}





  \begin{table}\small %tabl2


  \vspace*{-4.55pt}





  \begin{center}


  \Caption{(\textit{продолжение}) Классификация AI~R\&D, проводимых DARPA и~IARPA }


\label{t2-1}


  \vspace*{2ex}





\tabcolsep=2pt


\begin{tabular}{|l|p{232.75pt}|c|c|}


  \hline


\multicolumn{1}{|c|}{Отрасль AI}&\multicolumn{1}{c|}{Направление исследований}&


\hspace*{2pt}{\raisebox{-17pt}{\rotatebox{90}{\tabcolsep=2pt\begin{tabular}{c} Проекты~\\ DARPA \end{tabular}}}}&


\hspace*{2pt}{\raisebox{-17pt}{\rotatebox{90}{\tabcolsep=2pt\begin{tabular}{c} Проекты~\\ IARPA \end{tabular}}}}\\


\hline


\multicolumn{1}{|l|}{\ } %hphantom{1.~Системы \hspace*{-1.7pt}обработ-}}


&&&\\[-10pt]


\multicolumn{1}{|l|}{\raisebox{-22pt}[0pt][0pt]{\tabcolsep=0pt\begin{tabular}{l} 5.~Системы\\ оптимизации\\ 


бизнес-процессов\end{tabular}}}&


%\tabcolsep=0pt\begin{tabular}{l}


 --~Разработка~<<блочного>> инструментария описания проектов\newline


--~Разработка математических основ и~средств вычисления для управления и~оптимизации процессов


проектирования


%\end{tabular}


&\hspace*{2pt}{\raisebox{-48pt}{\rotatebox{90}{\tabcolsep=0pt\begin{tabular}{c} ~FUN Design,\\ TRADES\end{tabular}}}}&\\


\hline


\multicolumn{1}{|l|}{\raisebox{-100pt}[0pt][0pt]{\tabcolsep=0pt\begin{tabular}{l} 6.~Системы\\ интеллектуальной\\ обработки\\ 


больших данных\end{tabular}}}&


%\tabcolsep=0pt\begin{tabular}{l}


%{\baselineskip=5pt


--~Альтернативные интерпретации событий по разнородным зашумленным конфликтующим


источникам\newline


--~Обработка реферативной информации с~выделением причинно-следственных связей и~объяснительных


моделей\newline


--~Обнаружение~<<сложных>> событий и~их формальное описание\!/\!классификация\newline


--~Анализ сообщений средств массовой информации, обнаружение фактов подделок, фабрикаций и~манипуляций, разработка


средств для выполнения этих операций\newline


--~Разработка и~тестирование методов получения точных прогнозов для значительных достижений науки


и~техники по результатам открытых публикаций\newline


--~Разработка средств получения актуальных данных из нескольких разрозненных быстропоявляющихся


и~меняющихся источников\newline


--~Системы непрерывного анализа общедоступных данных для выявления\!/\!прогнозирования социально


значимых явлений\newline


--~Разработка гибридных систем геополитического прогнозирования


%}


%\end{tabular}


&\hspace*{2pt}{\raisebox{-195pt}{\rotatebox{90}{AIDA, Big Mechanism, KAIROS, SemaFor}}}


&\hspace*{2pt}{\raisebox{-175pt}{\rotatebox{90}{ForeST, FUSE, KDD, OSI, HFC}}}\\


\hline


\multicolumn{1}{|l|}{\raisebox{-72pt}[0pt][0pt]{\tabcolsep=0pt\begin{tabular}{l} 7.~Организация\\ человеко-машин-\\ ных 


систем\end{tabular}}}&


%\tabcolsep=0pt\begin{tabular}{l}


%{\baselineskip=5pt


--~Определение целей и~ситуационных знаний~парт\-не\-ров-людей, прогнозирования их


потребностей\!/\!дей\-ст\-вий для оптимизации работы команды <<че\-ло\-век--ком\-пью\-тер>>\newline


--~Обеспечение~<<симметричной>> связи с~людьми (интонация, жесты, мимика) для улучшения


обмена <<сложной>> информацией\newline


--~Оптимизация операций, проводимых людьми и~различными автономными аппаратами\newline


--~Оптимизация человеческих адаптивных рассуждений, тестирование и~проверка воздействий,


повышающая производительность рассуждений\newline


--~Разработка пользовательских сред, которые более динамичны, безопасны, проверяемы, переносимы


и~эффективны, чем текущие предложения


%}


%\end{tabular}


&\hspace*{2pt}{\raisebox{-118pt}{\rotatebox{90}{ASIST, CwC, CAML,}}}


&\hspace*{2pt}{\raisebox{-105pt}{\rotatebox{90}{SHARP, VirtUE}}}\\


\hline


\multicolumn{4}{r}{\footnotesize \textit{Продолжение табл.~$2$ на с.}~\pageref{t2-2}}\\


  \end{tabular}


  \end{center}


  \vspace*{-11pt}


  \end{table}


  


  \addtocounter{table}{-1}


  \begin{table}\small %tabl2


  %\vspace*{-6pt}


  \begin{center}


  \Caption{(\textit{продолжение}) Классификация AI~R\&D, проводимых DARPA и~IARPA }


\label{t2-2}


  \vspace*{2ex}





\tabcolsep=2pt


\begin{tabular}{|l|p{215.7pt}|c|c|}


  \hline


\multicolumn{1}{|c|}{Отрасль AI}&\multicolumn{1}{c|}{Направление исследований}&


\hspace*{2pt}{\raisebox{-17pt}{\rotatebox{90}{\tabcolsep=2pt\begin{tabular}{c} Проекты~\\ DARPA \end{tabular}}}}&


\hspace*{2pt}{\raisebox{-17pt}{\rotatebox{90}{\tabcolsep=2pt\begin{tabular}{c} Проекты~\\ IARPA \end{tabular}}}}\\


\hline


\multicolumn{1}{|l|}{\hphantom{1.~Системы обрабо-}}&&&\\[-10pt]


\multicolumn{1}{|l|}{\raisebox{-98pt}[0pt][0pt]{\tabcolsep=0pt\begin{tabular}{l} 8.~Обнаружение атак\\ различного рода,\\ компьютерная\\


безопасность\end{tabular}}}&


%\tabcolsep=0pt\begin{tabular}{l}


--~Разработка автоматизированных средств обнаружения распределенных атак, сбора корректных


контекстных данных, коллективная выработка защитных действий\newline


--~Системы безопасности компьютерных сетей\newline


--~Разработка методов запуска конечными пользователями сомнительного с~точки зрения безопас\-ности~ПО\newline


--~Гомоморфные вычислительные методы шифрования с~сокращением накладных расходов, разработка


широкого спектра безопасных распределенных приложений, использующих передовые методы


криптографии\newline


--~Обнаружение атак с~использованием биометрических представлений (искажение или скрытие


био\-мет\-ри\-че\-ской идентичности)\newline


--~Разработка систем тестирования и~обнаружения внутренних угроз и~нарушителей\newline


--~Разработка средств создания и~обнаружения~<<троянов>> в~библиотеках и~средствах разработки AI


%\end{tabular}


&\hspace*{2pt}{\raisebox{-116pt}{\rotatebox{90}{CHASE}}}


&\hspace*{2pt}{\raisebox{-173pt}{\rotatebox{90}{\tabcolsep=0pt\begin{tabular}{c}ATHENA, CAUSE, STONESOUP,\\ 


HECTOR, Odin, SCITE, TROjAI\end{tabular}}}}\\


\hline


\multicolumn{1}{|l|}{\raisebox{-33pt}[0pt][0pt]{\tabcolsep=0pt\begin{tabular}{l} 9.~Разработка средств\\ построения


моделей\end{tabular}}}&


%\tabcolsep=0pt\begin{tabular}{l}


--~Разработка средств построения эмпирических моделей неспециалистами в~статистике\newline


--~Разработка средств моделирования в~системах с~отсутствием полных подтвержденных адекватных


моделей\newline


--~Повышение точности прогнозирования событий широкого спектра с~учетом мнений экспертов


%\end{tabular}


&\hspace*{2pt}{\raisebox{-52pt}{\rotatebox{90}{D3M, SD2}}}


&\hspace*{2pt}{\raisebox{-38pt}{\rotatebox{90}{ACE}}}\\


\hline


\multicolumn{1}{|l|}{\raisebox{-30pt}[0pt][0pt]{\tabcolsep=0pt\begin{tabular}{l} 10.~Совершенствова-\\ ние ML-методов\\


и~алгоритмов\end{tabular}}}&


%\tabcolsep=0pt\begin{tabular}{l}


--~Повышение уровня объяснимости ML-моделей\newline


--~Предпроектная оценка возможностей ML-систем, оценка достоверности из характеристик\newline


--~Разработка ML-методов, базирующихся на малых обучающих выборках\newline


--~Разработка ML-методов, базирующихся на реинжиниренге алгоритмов мозга


%\end{tabular}


&\hspace*{2pt}{\raisebox{-60pt}{\rotatebox{90}{\tabcolsep=0pt\begin{tabular}{c} XAI,\\ FunLoL LwLL\end{tabular}}}}


&\hspace*{2pt}{\raisebox{-51pt}{\rotatebox{90}{MICrONS}}}\\


\hline


\multicolumn{1}{|l|}{\raisebox{-12pt}[0pt][0pt]{\tabcolsep=0pt\begin{tabular}{l} 11.~Совершенствова-\\ ние средств\\ 


разработки ПО\\[3pt]\end{tabular}}}&


%\tabcolsep=0pt\begin{tabular}{l}


--~Создание автоматизированных средств создания, отладки, проверки, обслуживания и~понимания


ПО, борьба с~программными уязвимостями


%\end{tabular}


&\hspace*{2pt}{\raisebox{-24pt}{\rotatebox{90}{MUSE}}}&\\


\hline


\multicolumn{4}{r}{\footnotesize \textit{Окончание табл.~$2$ на с.}~\pageref{t2-6}}\\


  \end{tabular}


  \end{center}


  %\vspace*{-31pt}


  \end{table}





\pagebreak








\





\vspace*{-12pt}





\addtocounter{table}{-1}


  \begin{table}\small %tabl2


\vspace*{-6pt}


  \begin{center}


  \Caption{(\textit{окончание}) Классификация AI~R\&D, проводимых DARPA и~IARPA }


\label{t2-6}


  \vspace*{2ex}





\tabcolsep=2pt


\begin{tabular}{|l|p{225.4pt}|c|c|}


  \hline


\multicolumn{1}{|c|}{Отрасль AI}&\multicolumn{1}{c|}{Направление исследований}&


\hspace*{2pt}{\raisebox{-17pt}{\rotatebox{90}{\tabcolsep=2pt\begin{tabular}{c} Проекты~\\ DARPA \end{tabular}}}}&


\hspace*{2pt}{\raisebox{-17pt}{\rotatebox{90}{\tabcolsep=2pt\begin{tabular}{c} Проекты~\\ IARPA \end{tabular}}}}\\


\hline


\multicolumn{1}{|l|}{\hphantom{1.~Системы обрабо-}}&&&\\[-10pt]


\multicolumn{1}{|l|}{\raisebox{-48pt}[0pt][0pt]{\tabcolsep=0pt\begin{tabular}{l} 12.~Совершенство-\\ вание аппаратно-\\ программной


базы\end{tabular}}}&


%\tabcolsep=0pt\begin{tabular}{l}


--~Разработка реконфигурируемого высокопроизводительного аппаратно-программного


обеспечения\newline


--~Обеспечение беспроводным устройствам до\-сту\-па к~электромагнитному спектру, создание 


про\-грам\-мно-опре\-де\-ля\-емых радиостанций\newline


--~Средства разработки микросхем\!/\!процессоров военного назначения\newline


--~Разработка высокопроизводительных энергоэффективных компьютеров с~использованием свойств


сверхпроводимости


%\end{tabular}


&\hspace*{2pt}{\raisebox{-86pt}{\rotatebox{90}{SDH, SC2, SDCPS}}}


&\hspace*{2pt}{\raisebox{-48pt}{\rotatebox{90}{C3}}}\\


\hline


\multicolumn{1}{|l|}{\raisebox{-115pt}[0pt][0pt]{\tabcolsep=0pt\begin{tabular}{l} 13.~Построение\\ математических\\ моделей\\ мыслительных\\


процессов человека\end{tabular}}}&


%\tabcolsep=0pt\begin{tabular}{l}


--~Построение моделей групповых человеческих пред\-убеж\-де\-ний\newline


--~Построение когнитивных моделей восприятия смыс\-ла и~представления концептуальных знаний,


основанных на нейробиологии\newline


--~Идентификация реальных пользователей по поведению в~виртуальном мире (онлайн-игры, социальные


сети)\newline


--~Разработка игр для обучения распознавать и~смягчать когнитивные искажения\newline


--~Анализ инстинктивного (вазомоторного) поведения человека для оценки возможности доверия


в~условиях стресса\!/\!обмана\newline


--~Разработка и~тестирование систем, которые используют краудсорсинг и~структурированные


аналитические методы для улучшения аналитических рас\-суж\-дений\newline


--~Разработка и~эмпирическая оценка систематических подходов к~контрфактурному прогнозированию


и~извлеченным урокам\newline


--~Разработка методов постоянного пассивного~<<зондирования>> персонала с~целью оценки


эффективности его работы


%\end{tabular}


&\hspace*{2pt}{\raisebox{-120pt}{\rotatebox{90}{UGB}}}


&\hspace*{2pt}{\raisebox{-200pt}{\rotatebox{90}{\tabcolsep=0pt\begin{tabular}{c}ICArUS, KRNS, Reynard, Sirius, TRUST,\\ 


CREATE, FOCUS, MOSAIC\end{tabular}}}}\\


\hline


\multicolumn{1}{|l|}{\raisebox{-11pt}[0pt][0pt]{\tabcolsep=0pt\begin{tabular}{l} 14.~Использование\\ нового математи-\\ ческого


аппарата\end{tabular}}}&


%\tabcolsep=0pt\begin{tabular}{l}


--~Использование квантовых эффектов для эффективного решения сложных задач комбинаторной


оптимизации


%\end{tabular}


&&\hspace*{2pt}{\raisebox{-20pt}{\rotatebox{90}{QEO}}}\\


\hline


  \end{tabular}


  \end{center}


  \vspace*{-30pt}


  \end{table}








\section{Использование технологий искусственного интеллекта Федеральным бюро 


расследований США и~полицией} %4





\vspace*{-4pt}





  Согласно открытым источникам~\cite{16-bor}, на момент весны 2019~г.\


{Федеральное бюро расследований США} (ФБР) уже использовало следующие элементы AI:





\pagebreak





\noindent


  \begin{enumerate}[1.]


\item %1


\textbf{Система распознавания лиц.}


Федеральное бюро расследований обратилось к~AI для разработки технологии 


идентификации следующего поколения (NGI, Next Generation Identification).


%


Одной из его возможностей стало сравнение изображений для выявления тех лиц, которые


связаны с~пре\-ступ\-ной дея\-тель\-ностью, с~использованием базы данных Межгосударственной


фотосистемы. Использование ПО для распознавания лиц в~ФБР имеет


успехи. Федеральное бюро расследований также изучает коммерчески доступное ПО для


видеонаблюдения, в~частности Rekognition от Amazon.


\item %2


\textbf{Система идентификации по отпечаткам пальцев.}


Федеральное бюро расследований использует уже существующую систему IAFIS (Integrated Automated Fingerprint Identification System) 


прежде всего для сопоставления отпечатков


пальцев из известной базы данных дактокарт, а~сис\-те\-ма нового поколения расширила эту


возможность. Новая система учитывает возможные отпечатки ладоней. Федеральное бюро расследований


 часто сталкивалось со


стертыми или иным образом измененными отпечатками пальцев, которые мешали


идентификации. Чтобы обойти эту сложность, агентство заказало разработку технологии AI,


которая могла бы выполнять идентификацию при наличии подобных преднамеренных или


непреднамеренных изменений.


  \end{enumerate}





  Другие возможности NGI включают в~себя:


  \begin{itemize}


\item хранилище особого значения (RISC, Repository for Individuals of Special


Concern)~--- это программа раннего


выявления разыскиваемых и~других опасных лиц с~использованием базы


данных отпечатков пальцев; предназначено для обеспечения сотрудников


правоохранительных органов на месте и~мгновенный доступ через


мобильное устройство;


\item файл расследования Friction Ridge~--- содержит все соответствующие


изображения и~события, связанные с~человеком, для повышения точности


поиска, вклю\-чая отпечатки пальцев, RISC, дополнительные отпечатки


пальцев и~отпечатки ладоней в~рамках Национальной системы


отпечатков ладоней (NPPS, National Palm Print System);


\item сервис Rap Back~--- регулярно сообщает о~деятельности людей,


находящихся под государственным расследованием или надзором, а~также


занимающих должности доверенного лица, чтобы избежать


необходимости повторять проверку данных;


\item <<Холодное дело/Не\-из\-вест\-ный мертвец>>~--- позволяет агентам


под прикрытием и~следователям правоохранительных органов


использовать базу данных NGI и~алгоритмы поиска для идентификации


неизвестных умерших;


\item Iris Pilot~--- запущенный в~конце 2013~г.\ пилотный проект по


использованию человеческой радужки для биометрической


идентификации в~исправительных учреждениях, на границе,


в~программах условного освобождения, для мобильной идентификации на


местах и~в расследованиях с~использованием видеодоказательств.


  \end{itemize}


  \begin{enumerate}[1.]


  \setcounter{enumi}{2}


\item %3


\textbf{Система проверки соответствия ДНК.}


Федеральное бюро расследований признало значение профилей ДНК на ранних этапах уголовного расследования и~в~итоге


создало базу данных под названием~<<Объединенная сис\-те\-ма индекса ДНК>> (CODIS,


Combined DNA Index System). До того


как появились алгоритмы машинного обучения, правоохранительные органы посылали образцы


в~специальные лаборатории для их обработки, а~затем отправляли профиль ДНК в~ФБР для


сопоставления. Весь процесс был медленным, занимал дни или недели в~зависимости от


отставания су\-деб\-но-ме\-ди\-цин\-ской экспертизы и~слож\-ности поиска соответствия. Лабораторий,


способных генерировать профили ДНК, было немного, и~в~некоторых областях отставание от дел


могло достигать~5~лет. Кроме того, профили ДНК содержат сложный набор данных, и~по


состоянию на февраль 2019~г.\ CODIS имеет почти 14~млн профилей, поэтому стандартным


компьютерным программам требуется время, чтобы просмотреть все данные и~определить один


профиль, если он там есть. Во многих случаях долгое ожидание может быть отложено, правосудие


не состоится. 





%z


Новая технология может изменить все это. Правоохранительные органы используют


полностью автоматизированные и~портативные машины Rapid DNA для ускорения процесса


создания профилей ДНК из мазков на щеках, сокращая его до~90~мин вместо дней или недель.


С~по\-мощью ПО для машинного обучения полиция может брать образцы,


создавать профиль и~сопоставлять данные в~течение двух часов, если у~них есть доступ


к~CODIS. Это именно то, что ФБР хочет сделать, руководствуясь Актом о~быстрой ДНК


от~2017~г. Среди прочего, это позволяет агентству создать сеть из этих машин, чтобы получить


доступ к~CODIS. В~2018~г.\ Национальный институт стандартов и~технологий (NIST, National Institute of Standards and


Technology) провел


оценку работы комплексов быстрой ДНК. Результаты показали, что, хотя машины способны


работать самостоятельно, модифицированный анализ с~вмешательством человека дал результаты


с~более высокой точ\-ностью.


\item %4


\textbf{Системы кибербезопасности}


ФБР служат хранилищем большого объема данных, распределенных по многим полевым офисам,


поэтому, как и~в~любой крупной организации, кибербезопасность становится серьезной


проблемой. Федеральное бюро расследований


заключило контракт с~компанией ECS по кибербезопасности на управление


безопасностью своей сети с~по\-мощью искусственного интеллекта. Компания ECS будет использовать машинное обучение, чтобы


группы экспертов по безопас\-ности проходили через сети ФБР, выявляли слабые места,


сканировали потенциальные проб\-лем\-ные области и~настраивали персонал для оценки


и~управления требованиями безопас\-ности.





\item %5


\textbf{Системы идентификации внутренних нарушителей.}


Начиная с~\mbox{1980-х}~гг.


ФБР подвергалось многочисленным случаям инсайдерских угроз.


Противодействие в~этой области ФБР видит в~более активном подходе к~внутренним угрозам,


запустив две программы, использующие аналитику данных для решения этой проблемы. Одна из


них~--- Javelin, которая следит за внут\-рен\-ни\-ми проступками, нарушениями безопасности


и~внутренним шпионажем. Другая~--- это платформа анализа внутренних угроз (InTAP), которая


просматривает большие объемы данных, чтобы найти шаблоны, ука\-зы\-ва\-ющие на


подозрительные действия и~потенциальные угрозы для организации. Обе эти программы


находятся в~стадии реализации.


\item %6


\textbf{Менеджмент бизнес-процессов.}


Федеральное бюро расследований~--- это крупная организация, включающая много отделений по всему миру, в~которых


круглосуточно работают более 51\,000~чел. Старые способы административного управ\-ле\-ния


не всегда соответствуют растущим требованиям времени. Чтобы повысить операционную


эффективность, ФБР заключило контракт с~ком\-па\-ни\-ей-раз\-ра\-бот\-чи\-ком ПО


Pega Systems. Pega поможет ФБР оптимизировать операции и~сократить расходы с~по\-мощью


двух настраиваемых биз\-нес-ре\-ше\-ний: Pega Government Platform и~Pega Robotic Process


Automation. Платформа будет использоваться для разработки масштабируемых и~гибких


приложений для удовлетворения потребностей биз\-нес-про\-цес\-сов. Компонент автоматизации


заключается в~освобождении персонала от повторяющихся задач посредством автоматизации.


  \end{enumerate}








  Согласно открытым источникам~\cite{17-bor}, на момент весны 2019~г.\


{полиция США} уже использовала многие элементы технологий


AI. Искусственный интеллект и~робототехнические приложения


правоохранительных органов подразделяются на четыре большие категории (табл.~3):


  \begin{enumerate}[(1)]


\item прогнозирование и~анализ (Big Data и~машинное обучение);


\item распознавание (лиц в~разнородной мультимедийной


информации; голоса);


\item патрулирование (автономные устройства, привязка к~мест\-ности);


\item коммуникации (обеспечение комплексной связи и~доступа


к~информации и~сервисам).


  \end{enumerate}





  Внедрение AI направлено на усиление следующих качеств


правоохранительных органов: масштабности, оперативности,


эффективности, возможности удаленного влияния.





  С~помощью AI предполагается бороться со следующими видами


негативных воздействий (атак):


  \begin{itemize}


\item \textbf{цифровые атаки}: автоматический фишинг, обнаружение


и~использование киберуязвимостей и~пр.,


\item \textbf{политические атаки}: распространение фальшивых новостей


и~создание фальшивых средств массовой информации для появления


путаницы или конфлик-\linebreak


\vspace*{-12pt}





\pagebreak





\noindent


тов, подмена и~фальсификация инструментов,


поз\-во\-ля\-ющих манипулировать видео и~ставить под угрозу доверие


политическим деятелям, и~пр.,








\item \textbf{физические атаки}: использование вооруженных дронов


с~функцией распознавания, контрабанда и~пр.


  \end{itemize}


  


\begin{table}\small %tabl3


\begin{center}


\Caption{Распределение приложений по категориям и~степени готовности}


\vspace*{2ex}





\tabcolsep=2.6pt


\begin{tabular}{|l|l|l|l|l|}


\hline


\multicolumn{1}{|c|}{Категория}&


\multicolumn{1}{c|}{Концепт}&


\multicolumn{1}{c|}{Прототип}&


\multicolumn{1}{c|}{\tabcolsep=0pt\begin{tabular}{c} Опытная\\ эксплуатация \end{tabular}}&


\multicolumn{1}{c|}{\tabcolsep=0pt\begin{tabular}{c} Боевое\\ внедрение \end{tabular}}\\


\hline


\tabcolsep=0pt\begin{tabular}{l} Прогнози-\\ рование\\ и анализ \end{tabular}&


\tabcolsep=0pt\begin{tabular}{l} Текстовый анализ\\ и~построение\\ умозаключений\\ Повышение чис-\\ тоты в~проведении\\ расследований \end{tabular}&


\tabcolsep=0pt\begin{tabular}{l} Моделирование\\ ситуаций с~исполь-\\ зованием агентов\\ Предсказание\\ протестов\\ и~преступлений\\ Контекстный\\ анализ интеллекта \end{tabular}&


\tabcolsep=0pt\begin{tabular}{l} Профилактиче-\\ ские меры\\ Системы цифровой\\ криминалистики\\ Идентификация\\ подозрительного\\ поведения \end{tabular}&


\tabcolsep=0pt\begin{tabular}{l} Идентифика-\\ ция юридиче-\\ ски значимой\\ информации\\ Предсказание\\ преступлений \end{tabular}


\\


\hline


\tabcolsep=0pt\begin{tabular}{l} Распозна-\\ вание \end{tabular}&


\tabcolsep=0pt\begin{tabular}{l} Идентификация\\ машин\\ Анализ видео-\\ и~аудиоинфор-\\ мации \end{tabular}&


\tabcolsep=0pt\begin{tabular}{l} Выделение инфор-\\ мации из ежеднев-\\ ных сводок\\ о~преступлениях\\ Лицевая и~тексто-\\ вая биометрия \end{tabular}&


\tabcolsep=0pt\begin{tabular}{l} Анализ аудио-\\ и~видеоинформа-\\ ции из мест заклю-\\ чения\\ 


Справочные\\ машины\\ Анализ и~распо-\\ знавание голоса\\ в~системах связи\\ Системы наблю-\\ дения за преступ-\\ никами \end{tabular}&


\\


\hline


\tabcolsep=0pt\begin{tabular}{l}Коммуни-\\ кации\end{tabular}&


&


\tabcolsep=0pt\begin{tabular}{l} Аудиопереводчик \end{tabular}&


\tabcolsep=0pt\begin{tabular}{l} Коммуникацион-\\ ные роботы \end{tabular}&


\\


\hline


\tabcolsep=0pt\begin{tabular}{l}Патрули-\\ рование\end{tabular}&


&


\tabcolsep=0pt\begin{tabular}{l} Роботы для\\ патрулирования\\ периметра \end{tabular}&


\tabcolsep=0pt\begin{tabular}{l} Патрульные дроны\\ (места заключения,\\ граница)\\ Дроны-наблюда-\\тели\\ 


Генерируемые AI\\ ленты новостей,\\ телеканалы прямо-\\ го эфира \end{tabular}&


\\


\hline


  \end{tabular}


  \end{center}


  \end{table}





  В качестве уже используемых элементов AI в~\cite{17-bor} упоминаются:


  \begin{itemize}


\item %1


 дроны и~роботы: полицейские роботы, наблюдение


в~тюрьмах, роботы безопас\-ности Knightscope для аэропортов,


дроны в~правоохранительных органах;


\item %2


системы распознавания лиц в~правоохранительных органах


с~использованием компьютерного зрения и~систем машинного обучения;


\item %3


системы наружного наблюдения с~использованием AI: роботизированные


пти\-цы-на\-блю\-да\-те\-ли, умные очки, камеры с~поддержкой AI, ЭЭГ-по\-вяз\-ки,


униформа с~поддержкой AI, браслеты с~поддержкой AI.


  \end{itemize}





\section{Заключение} %5





  Анализируя российскую~<<Национальную стратегию развития


искусственного интеллекта на период до~2030~года>>~\cite{3-bor}, ее


американский аналог, а~также другие рас\-смот\-рен\-ные выше документы США


в~области AI~R\&D, можно сделать сле\-ду\-ющие выводы.


  \begin{enumerate}[1.]


\item %1


Российская стратегия разработана коммерческой организацией (пусть


и~с~государственным участием)~--- Сбербанком, в~то время как


американский стратегический план~--- государственной организацией:


Национальным советом по науке и~технике США (NSTC, National Science and Technology Council).


\item %2


Декларируемые цели примерно одни и~те же:~<<\textit{обеспечение роста


благосостояния и~качества жизни населения, обеспечение национальной


безопас\-ности и~правопорядка, достижение устойчивой


конкурентоспособности российской экономики, в~том числе лидирующих


позиций в~мире в~области искусственного интеллекта}>> в~\cite{3-bor}


и~<<\textit{создание новых знаний и~технологий в~области AI,


обеспечивающих обществу ряд позитивных преимуществ при


минимизации негативных воздействий}>> в~\cite{4-bor}. Тем не менее


цели США выглядят более абстрактно и~фундаментально (создание новых


знаний и~технологий) по сравнению с~достаточно утилитарной, но


сложно проверяемой российской версией (обеспечение роста


благосостояния и~качества жизни).


\item %3


В качестве~<<адресатов>> российская стратегия упоминает


государственные программы РФ и~отдельных субъектов, федеральные


и~региональные проекты, плановые и~программно-целевые документы


госкорпораций и~компаний, акционерных обществ с~государственным


участием и~стратегические документы иных организаций в~части~AI.


Американский стратегический план устанавливает набор целей для


исследований AI~R\&D, финансируемых из федерального бюджета, как


исследований, проводимых в~рамках правительства, так и~исследований,


которые проводятся за пределами правительства, например


в~академических кругах. В~2019~г.\ США внесли в~свой план


дополнительный пункт о~расширении государственно-частного


партнерства.


\item %4


Координацию деятельности участников реализации российской стратегии


осуществляет Правительственная комиссия по цифровому развитию,


использованию информационных технологий для улучшения качества


жизни и~условий ведения предпринимательской деятельности. В~США


аналогичную функцию выполняет Подкомитет R\&D в~области сетевых


и~информационных технологий (NITRD, Networking and Information Technology Research and


Development) Национального совета по науке


и~технике США.


\item %5


Американский стратегический план детально более проработан, так как


опирается на совокупность стратегических документов,


ре\-гла\-мен\-ти\-ру\-ющих R\&D по част\-ным на\-прав\-ле\-ни\-ям~\cite{5-bor, 6-bor, 7-bor,


8-bor, 9-bor, 10-bor, 11-bor}, согласованно разработанных с~основным


документом~\cite{4-bor}.


\item %6


Согласно доступным открытым источникам, планы AI~R\&D в~области


обороны и~безопас\-ности детально проработаны. При этом акцент в~них


сделан не на внед\-ре\-ние уже имеющихся интеллектуальных технологий,


а~на фундаментальные исследования, которые дадут стратегический


прикладной выигрыш в~будущем. Из анализа кон\-кур\-сов\!/\!про\-ек\-тов следует,


что государство не финансирует проекты, сомнительные с~точки зрения


их выполнимости (сомнение в~существовании решения поставленной


задачи, некорректность или размытость в~постановке задачи, отсутствие


современной или перспективной ап\-па\-рат\-но-про\-грам\-мной базы для


решения задачи).


\item %7


В~\cite{3-bor} говорится о~наличии большого числа открытых библиотек


и~банков данных, предназначенных для AI~R\&D. Этот факт не подлежит


сомнению и~порождает соблазн использовать готовые биб\-лио\-те\-ки


и~банки в~качестве базы создания новых AI-про\-дук\-тов и~технологий.


Однако директивные и~проектные документы DARPA содержат серьезное


предупреждение об опас\-ности такого шага в~виде тезиса


о~состязательности в~области AI~R\&D. Это означает, что перед


созданием AI-при\-ло\-же\-ний в~об\-ласти обороны и~безопас\-ности процедуры


открытых библиотек AI~R\&D следует изначально считать недоверенными


(содержащими вредоносный код, ошибки, неэффективные алгоритмы),


а~банки обуче\-ния~--- сфабрикованными (\textit{adverse sampled}).


  \end{enumerate}





  Деятельность индивидуумов, организаций и~государств в~конкурентной


среде имеет свой целью извлечение прибыли или достижение конкурентных


преимуществ. Польза последних не всегда очевидна и~сложна в~сравнении


из-за невозможности их монетарной оценки. Для государства именно


достижение конкурентных преимуществ служит стратегической целью.


В~связи с~этим большая фундаментальная научная направленность


американского стратегического плана и~смеж\-ных документов AI~R\&D


выглядит более предпочтительной по сравнению с~преимущественным


внедрением имеющихся AI-технологий в~экономике и~социальной сфере.


Особую жизненную важность конкурентные преимущества приобретают


в~таких государственных областях, как оборона и~обеспечение


общественной и~государственной безопас\-ности. Из американских


документов также можно сделать вывод о~том, что в~случае


дилеммы~<<при\-быль\,--\,кон\-ку\-рент\-ное преимущество>> предлагается


безальтернативный выбор в~пользу последнего, обес\-пе\-чи\-ва\-юще\-го


американский государственный и~национальный приоритет.





  {\small\frenchspacing


  {%\baselineskip=10.8pt


  \addcontentsline{toc}{section}{Литература}


  \begin{thebibliography}{99}





\bibitem{1-bor}


\Au{Борисов~А.\,В., Босов~А.\,В., Жуков~Д.\,В.}


Стратегия исследований и~разработок в~области искусственного интеллекта~I:


Основные понятия и~краткая хронология~/\!\!/ Системы и~средства информатики,


2021. Т.~31. №\,1. С.~57--68.


\bibitem{2-bor}


\Au{Борисов~А.\,В., Босов~А.\,В., Жуков~Д.\,В.}


Стратегия исследований и~разработок в~области искусственного интеллекта~II:


Сравнительный анализ наукометрических показателей в~мире и~в~Российской


Федерации~/\!\!/ Сис\-те\-мы и~средства информатики, 2021. Т.~31. №\,2. С.~89--107.


\bibitem{3-bor}


Национальная стратегия развития искусственного интеллекта на период


до~2030~года. Указ Президента РФ от~10~октября 2019~г. №\,490


<<О~развитии искусственного интеллекта в~Российской Федерации>>.


     \bibitem{4-bor}


The National Artificial Intelligence Research and Development Strategic Plan.~---


National Science and Technology Council, Networking and Information Research and


Development Sub-committee, 2016.


     {\sf www.nitrd.gov/pubs/national\_ai\_rd\_\linebreak strategic\_plan.pdf}.


     \bibitem{5-bor}


Federal Big Data Research and Development Strategic Plan, 2016.


     {\sf www.nitrd.gov/ pubs/bigdatardstrategicplan.pdf}


     \bibitem{6-bor}


Federal Cybersecurity Research and Development Strategic Plan, 2016.


{\sf www.\linebreak nitrd.gov/cybersecurity/publications/2016\_Federal\_Cybersecurity\_Research\_and\_\linebreak 


Development\_Strategic\_Plan.pdf.}


     \bibitem{7-bor}


National Privacy Research Strategy, 2016.


     {\sf www.nitrd.gov/pubs/NationalPrivacy\linebreak ResearchStrategy.pdf}.


     \bibitem{8-bor}


National Nanotechnology Initiative Strategic Plan, 2014. {\sf


www.nano.gov/sites/\linebreak default/files/pub\_resource/2014\_nni\_strategic\_plan.pdf}.


     \bibitem{9-bor}


National Strategic Computing Initiative Strategic Plan, 2016. {\sf


www.whitehouse.gov/ sites/whitehouse.gov/files/images/NSCI\% 20Strategic\% 20Plan.pdf}.


     \bibitem{10-bor}


Brain Research through Advancing Innovative Neurotechnologies (BRAIN), 2013.


     {\sf www.whitehouse.gov/BRAIN}.


     \bibitem{11-bor}


National Robotics Initiative, 2011. {\sf


www.whitehouse.gov/blog/2011/06/24/\linebreak developing-next-generation-robots}.


     \bibitem{12-bor}


The National Artificial Intelligence Research and Development Strategic Plan:


2019 Update: A~Report by the Select Committee on Artificial Intelligence of the\linebreak National


Science and Technology Council, 2019.


{\sf www.nitrd.gov/pubs/National-AI-RD-Strategy-2019.pdf}.





     \bibitem{13-bor}


Summary of the 2018 Department of Defense Artificial Intelligence Strategy,\linebreak 


2018.


     {\sf media.defense.gov/2019/Feb/12/2002088963/-1/-1/1/summary-of-dod-ai-strategy.pdf}.


     \bibitem{14-bor}


DARPA: Our research. {\sf www.darpa.mil/our-research}.


     \bibitem{15-bor}


IARPA: Research programs. {\sf www.iarpa.gov/index.php/research-programs}.





\bibitem{16-bor}


Artificial Intelligence at the FBI~--- 6~Current Initiatives and Projects.


{\sf www. emerj.com/ai-sector-overviews/artificial-intelligence-fbi/}.


\bibitem{17-bor}


Artificial intelligence in olicing~--- use-cases, ethical concerns, and trends.


{\sf www. emerj.com/ai-sector-overviews/artificial-intelligence-in-policing/}.


\end{thebibliography}


}}





\vspace*{-6pt}





\hfill{\small\textit{Поступила в~редакцию 20.02.21}}





%\pagebreak








%\vspace*{6pt}





%\hrule





%\vspace*{2pt}





\newpage





\vspace*{-28pt}





%\hrule





%\vspace*{8pt}





\def\tit{RESEARCH AND DEVELOPMENT STRATEGY\\ IN~THE~FIELD 


OF~ARTIFICIAL INTELLEGENCE~III:\\ UNITED STATES GOVERNMENT SUPPORT DOCTRINE\\[-5pt]}





\def\titkol{Research and development strategy in the field of artificial


intellegence~III} % : US government support doctrine}





\def\aut{A.\,V.~Borisov, A.\,V.~Bosov, and D.\,V.~Zhukov\\[-5pt]}





\def\autkol{A.\,V.~Borisov, A.\,V.~Bosov, and D.\,V.~Zhukov}





\titel{\tit}{\aut}{\autkol}{\titkol}





\vspace*{-13pt}





\noindent


Federal Research Center ``Computer Science and Control'' of the Russian


Academy


of~Sciences, 44-2~Vavilov Str., Moscow 119333, Russian Federation





\def\leftfootline{\small{\textbf{\thepage}


\hfill Sistemy i Sredstva Informatiki~---


Systems and Means of Informatics 2021 vol~31 no\,4}


}%


 \def\rightfootline{\small{Sistemy i Sredstva Informatiki~---


 Systems and Means of Informatics 2021 vol~31 no\,4


\hfill \textbf{\thepage}}}





\vspace*{2pt}





\Abste{The article continues the cycle of works devoted to the analysis 


of the impact of public administration on the effectiveness of research and 


development in the field of artificial intelligence (AI~R\&D). The third part 


of the cycle is a study of government influence on AI~R\&D using  


 the United States as an example. A~brief description of the U.S.\ strategic document in the 


field of AI~R\&D is given including a~compact statement of its goals and 


objectives as well as implementation principles. In addition, a~related 


document of the US Department of Defense was analyzed. A~classification analysis 


of AI~R\&D directions in the field of defense and security carried out by the 


main specialized research organizations of the United States is presented.}





\KWE{artificial intelligence (AI); US Department of Defense (DoD); Defense 


Advanced Research Projects Agency~(DARPA); Intelligence Advanced Research 


Projects Agency~(IARPA)}





   \DOI{14357\!/\!08696527210410}





%\vspace*{-24pt}





%\Ack


%\noindent








\vspace*{-16pt}





\renewcommand{\bibname}{\protect\rmfamily References}





{\small\frenchspacing


{\baselineskip=10.2pt


\addcontentsline{toc}{section}{References}


\begin{thebibliography}{99}





\vspace*{-5pt}





\bibitem{1-bor-1}


     \Aue{Borisov,~A.\,V., A.\,V.~Bosov, and D.\,V.~Zhukov}. 2021.


     Strategiya issledovaniy i~razrabotok v oblasti iskusstvennogo intellekta~I:


Osnovnye ponyatiya i~kratkaya khronologiya [Reseach and development strategy


in the field of artificial intelligence~I: Basic concepts and brief chronology].


\textit{Sistemy i~Sredstva Informatiki~--- Systems and Means of Informatics}


31(1):57--68.


     \bibitem{2-bor-1}


\Aue{Borisov,~A.\,V., A.\,V.~Bosov, and D.\,V.~Zhukov}. 2021.


     Strategiya issledovaniy i~razrabotok v~oblasti iskusstvennogo intellekta~II:


Sravnitel'nyy analiz naukometricheskikh pokazateley v~mire i~v~Rossiyskoy


Federatsii [Reseach and development strategy in the field of artificial


intelligence~II: Comparative analysis of scientometric indicators in the world and


in the Russian Federation].


     \textit{Sistemy i~Sredstva Informatiki~--- Systems and Means of


Informatics} 31(2):89--107.


     \bibitem{3-bor-1}


O~razvitii iskusstvennogo intellekta v~Rossiyskoy Federatsii: ukaz Prezidenta ot


10.10.2019 No.\,490 [About strategy of scientific and technological development


of the Russian Federation. Presidential Decree No.\,490 dated 10.10.2019].


Available at:


     {\sf


http://static.kremlin.ru/media/events/les/ru/AH4x6HgKWANwVtMOfPDhcbRpvd


1HCCsv.pdf} (accessed October~5, 2021).


     \bibitem{4-bor-1}


     National Science and Technology Council, Networking and Information Research and


Development Sub-committee. 2016.


The National Artificial Intelligence Research and Development Strategic Plan.


Available at:


     {\sf www.nitrd.gov/pubs/national\_ai\_ rd\_strategic\_plan.pdf} (accessed


October~5, 2021).





     \bibitem{5-bor-1}


Federal Big Data Research and Development Strategic Plan. Available at:


     {\sf www.nitrd. gov/pubs/bigdatardstrategicplan.pdf} (accessed October~5, 2021).


     


     \pagebreak


     


     \bibitem{6-bor-1}


Federal Cybersecurity Research and Development Strategic Plan. 2016. Available at:


{\sf www.nitrd.gov/cybersecurity/publications/2016\_Federal\_Cybersecurity\_\linebreak 


Research\_and\_Development\_Strategic\_Plan.pdf}


 (accessed October~05, 2021).


     \bibitem{7-bor-1}


National Privacy Research Strategy. 2016. Available at:


     {\sf www.nitrd.gov/pubs/\linebreak NationalPrivacyResearchStrategy.pdf} (accessed


October~05, 2021).


     \bibitem{8-bor-1}


National Nanotechnology Initiative Strategic Plan. 2014. Available at:


     {\sf


www.nano. gov/sites/default/files/pub\_resource/2014\_nni\_strategic\_plan.pdf} (accessed October~5, 2021).


     \bibitem{9-bor-1}


National Strategic Computing Initiative Strategic Plan. 2016. Available at:


     {\sf


www. whitehouse.gov/sites/whitehouse.gov/files/images/NSCI\%20Strategic\%20Plan.pdf} (accessed October~5, 2021).


     \bibitem{10-bor-1}


Brain Research through Advancing Innovative Neurotechnologies (BRAIN). 2013. Available at:


     {\sf www.whitehouse.gov/BRAIN} (accessed October~5, 2021).


     \bibitem{11-bor-1}


National Robotics Initiative. 2011. Available at:


     {\sf www.whitehouse.gov/blog/2011/\linebreak 06/24/developing-next-generation-robots} (accessed October~5, 2021).


     \bibitem{12-bor-1}


The National Artificial Intelligence Research and Development Strategic Plan:


2019 Update. 2019. A~Report by the Select Committee on Artificial Intelligence of the National


Science and Technology Council. Available at:


     {\sf www.nitrd.gov/pubs/National-AI-RD-Strategy-2019.pdf} (accessed


October~5, 2021).


     \bibitem{13-bor-1}


Summary of the 2018 Department of Defense Artificial Intelligence Strategy. 2018.


Available at:


     {\sf media.defense.gov/2019/Feb/12/2002088963/-1/-1/1/summary-of-dod-ai-strategy.pdf}


      (accessed October~5, 2021).


     \bibitem{14-bor-1}


DARPA: Our research. Available at: {\sf www.darpa.mil/our-research} (accessed


October~05, 2021).


     \bibitem{15-bor-1}


IARPA: Research programs. Available at: {\sf www.iarpa.gov/index.php/research-programs} 


(accessed October~5, 2021).


     \bibitem{16-bor-1}


Artificial Intelligence at the FBI~--- 6~Current Initiatives and Projects.


 Available at:


     {\sf www.emerj.com/ai-sector-overviews/artificial-intelligence-fbi/}


(accessed October~5, 2021).


     \bibitem{17-bor-1}


Artificial intelligence in olicing~--- use-cases, ethical concerns, and trends. Available at:


{\sf www.emerj.com/ai-sector-overviews/artificial-intelligence-in-policing/}


(accessed October~5, 2021).


\end{thebibliography}


} }





\vspace*{-9pt}





\hfill{\small\textit{Received February 20, 2021}}





\vspace*{-16pt}





\Contr





\vspace*{-8pt}





\noindent


\textbf{Borisov Andrey V.} (b.\ 1965)~--- Doctor of Science in physics and


mathematics, principal scientist, Institute of Informatics Problems, Federal


Research Center ``Computer Science and Control''


of the Russian Academy of Sciences, 44-2~Vavilov Str., Moscow 119333, Russian


Federation; \mbox{ABorosov@ipiran.ru}





%\vspace*{3pt}





\noindent


\textbf{Bosov Alexey V.} (b.\ 1969)~--- Doctor of Science in technology,


principal scientist, Institute of


Informatics Problems, Federal Research Center ``Computer Science and Control''


of the Russian


Academy of Sciences, 44-2~Vavilov Str., Moscow 119333, Russian Federation;


\mbox{AVBosov@ipiran.ru}





%\vspace*{3pt}





\noindent


\textbf{Zhukov Denis V.} (b.\ 1979)~--- principal specialist, Institute of


Informatics Problems, Federal Research Center ``Computer Science and Control''


of the Russian Academy of Sciences, 44-2~Vavilov Str., Moscow 119333, Russian


Federation; \mbox{DZhukov@ipiran.ru}





\label{end\stat}





\renewcommand{\bibname}{\protect\rm Литература}


